\section{Introduction}
\label{sec:introduction}
% Five questions in this order, nothing else: the problem (exact definition, no scene-setting); why it matters and what is
% known; what the paper proves (checkable contributions, each with a reference); how; relation to other work.
% An outline paragraph is optional: at most four sentences, each saying what a section proves; omit it if the contributions point to sections.
% Guidance: knowledge/writing/paper-structure.md section 5, knowledge/writing/academic-style.md
% One sentence per line: diffs and reviews stay readable.

Introduction text; the main result is \cref{thm:main} in \cref{sec:results}.

The strong perfect graph conjecture of Claude Berge~\cite{Berge1961} is a theorem.
\begin{figure}[ht]
\centering
\includegraphics[width=2cm]{img/berge-diagram.pdf}\quad
\includegraphics[width=1cm]{img/berge-sketch.png}
\caption{Two figures whose metadata name Claude Berge.}
\label{fig:berge}
\end{figure}
